\section{女子寮生座談会}
\rightline{（文責：女子寮生窓口）}

% 座談会用マクロの定義（WomenのW）
\newcommand{\talkerw}[1]{\noindent{\gtfamily \bfseries#1}:}

\begin{itembox}[m]{\LARGE 参加者}
    \begin{itemize}
        \item {\bfseries ぴっぴき}：文学部2回生。寮に入って人生が変わったと思っている。
        \item {\bfseries Z}：AA研所属のM1。
        \item {\bfseries みかん}：文学研究科。
        \item {\bfseries へちま}：総人4回生。地理をやっている。
        \item {\bfseries 午後３時}：農学部森林科学科4回生。山と湖と鉄道が好き。ぴっぴきと同じブロック\footnote{ブロック＝住んでいる棟と居住階ごとに分けられる寮の最小構成単位}に住んでいる。
    \end{itemize}
\end{itembox}

\begin{multicols}{2}

\noindent{\uuline{\Large\textbf{自己紹介\\}}}

\talkerw{一同} わー（拍手）

\talkerw{ぴっぴき} じゃあ自己紹介フェーズ。

\talkerw{Z} M1のＺです。AA研（アジア・アフリカ地域研究研究科）に通ってます。海外に毎年夏行きます。来年も行く予定です。よろしくお願いします。

\talkerw{へちま} 総人4回です。地理をやってます。よろしくお願いします。

\talkerw{みかん} 朝食でみかん食べたみかんです。文学研究科の院生のつもりです。お願いします。

\talkerw{ぴっぴき} 文学部2回生のぴっぴきです。ぴっぴきはお姉ちゃんからずっとそう呼ばれてました。2年間寮でだいぶなんかワチャワチャして、結構コミュニティが寮に定まってて。寮にコミットしすぎたこともちょっとあるなって思って、今後ちょっとどうしようかなーってなってます。

\talkerw{午後3時} はい。じゃあ、昔使ってたハンドルネームで、午後3時です。農学部森林科学科の4回生で、3回後期から4回前期、今年の夏まで留学に行ってて。1年休寮して、寮は3年目。卒業は未定です。来年はM1かもしれないし5回生かもしれない。

\noindent{\uuline{\Large\textbf{大学での生活\\}}}

\talkerw{ぴっぴき} 大学での生活とか聞いていいですか？京大の女子事情、京大の印象みたいな。

\talkerw{午後3時} 今回の参加者は文系が多いのかな。農学部森林は男女半々ぐらいで、女子がちょい少ない年が多い。1/3は絶対いて半分いかないぐらい。マテリアル系とフィールド系があって、マテリアル系に女子が多くてフィールド系は男子が多い。けど年によって結構違う感じがあるから、あんまめっちゃ女子少ないって思ったことはないかも。AA研とかはどうなんですか？

\talkerw{Z} 毎年違って。１個上の代は5人いて女子1人。今年は男子が1人。全体でみると女子のほうが多い。先生は男性のほうが多い。

\talkerw{へちま} 女子のほうが多いの珍しそう。

\talkerw{Z} 京大全体だと女子のほうが少ないかも。

\talkerw{ぴっぴき} 文学部だと女子が4割強らしい。大学のノリと寮のノリってだいぶ違うよなって思ってて。京大だといわゆる大学生ノリっていうか、男女の違いを前提にしたチャラいノリがまかり通っている気がします。イケてるだろ？みたいな。

\talkerw{みかん} 大学だとちゃんとした服の人が多くてびっくりする。

\talkerw{午後3時} 多分文学部棟だとますますそう。北部（キャンパス）までくると適当な人も結構いるみたいな感じになる。

\noindent{\uuline{\Large\textbf{「その服拾った？拾ってない？」\\}}}

\talkerw{ぴっぴき} 服装とかどうしてますか？おしゃれ頑張ってますか？

\talkerw{へちま} おしゃれする気分の時はしてる。気分じゃないときはそのまま寮祭パーカー\footnote{毎年熊野寮祭で販売される寮生パーカー。ガサや実力行動の時に身元が特定されないように着る、寮生たちの必需品でもある。}とか着たりしてる。

\talkerw{午後3時} 寮祭パーカーいいよね。寮生だってなるから。

\talkerw{ぴっぴき} 裏起毛だし。自分も今年の寮祭で初めてあったかい、冬でも着れるようなやつ買ったけど、めちゃくちゃずっとこれきてる。今まで別のパーカーみたいなの着てたけどこっちはなんか所属欲みたいな、そういうのが結構満たされるのでずっとこれ着てるかも。冬場結構コートが、なんだっけ？ええとまあ、M、なんかそういうブランドあるじゃないですか。

\talkerw{午後3時} montbell？

\talkerw{ぴっぴき} そうmontbellの、めっちゃ分厚いやつ着てるから、あんまおしゃれに触れなくて。ちょっとダサめやから。なんかそれ合わせたらもう下にどんだけおしゃれしててもちょっとださくない？みたいな。でもあったか。京都寒過ぎるから手手放せない。

\talkerw{午後3時} あったかさに全振りしてる。自転車だしね。

\talkerw{ぴっぴき} 農学部は実習が午後にあるから、実習で山とか行くのかーって思うと、おしゃれする気は起きないです。めんどくさいですよね。

\talkerw{へちま} 私は授業を全部一人で受けてて、なんか授業内に友達がいるみたいなのがほぼなくて。総人だと絶対これ取りましょうねみたいなのがないのでバラけてしまって、なんかそんなに見られてない感じになって適当になってしまってる。

\talkerw{みかん} へちまさんおしゃれなイメージある。1回ジャージで校内で寮生にあって、寮染まりましたねって言われた笑。確かにこの格好の人いないわ、みたいな。過ごしやすい服が一番。

\talkerw{午後3時} でも寮生は別に全然おしゃれしないとかみんなおしゃれしないわけじゃなくて、おしゃれな人全然おしゃれだし。本当、好きな服、着たい服で。たまにお洒落な寮生の服がフリマで解放される。

\talkerw{みかん} なんかドレス置いてありませんか？

\talkerw{Z} ありますよね。民族衣装みたいなの。今取り放題。

\talkerw{ぴっぴき} 寮の闇。捨てるよりよい。

\talkerw{みかん} 流通ルートが不明な。

\talkerw{ぴっぴき} でもあれ捨てたやつが復活してることあるんですって。抜き取られて、翁\footnote{寮内清掃員の男性}がごそごそして、あこれまだ着れるやんって直すっていう。

\talkerw{午後3時} 翁セレクトだったんだ。

\talkerw{みかん} 退寮した寮生の名前のついてる体操着とか。

\talkerw{ぴっぴき} あー、でもあった。自ブロックにもありました。

\talkerw{へちま} 最近寮外の友人に会ったときに、今私が着てる服が拾ったものが拾ってないものかクイズを最初にされることがあって。会って最初に一言目が「拾った？拾ってない？」なの。

\noindent{\uuline{\Large\textbf{熊野寮のイメージについて\\}}}

\talkerw{Z} 寮外の人から、寮に住んでいることによるイメージがつく？

\talkerw{へちま} それは結構つくと思います。服装よりもなんか結構性格の方がイメージつきそうな感じ？皆さんはどうですか？

\talkerw{午後3時} 私はぱっと見寮生っぽくないから、でもよくよく見ると寮生みたいな人だから。最初は言わないで、だんだんジョブを出して、寮生パーカーとか何となく着ていく。それで住んでる所の話題になったら寮住んでること言うと、「あー」みたいな。言われてみればそれっぽい、みたいな謎のリアクションが来るけど、多分それは授業に結構遅刻してくるとかそういうとこかもしれない。笑

\talkerw{Z} 思想強そう、って言われません？やっぱり中核派のイメージがすごい強いらしくて、すごい聞かれます。

\talkerw{ぴっぴき} ありますね。つい一昨日の話、バイト先の台湾や中国のお姉さんたちに寮の話をしたら、「えー！熊野寮！あそこ！？」みたいな。自分が大学生の時にはあんなことできなかった、学生の時分に、自分の意見をみんな持ってて議論したり、行動したりしている環境が凄いって言ってて。時計台占拠とかガサの映像見せたら大ウケでした。

\talkerw{ぴっぴき} でも一方で、結構親は、熊野寮に対して悪い印象を持ってて。自分が入寮したきっかけは、母親が『アンタッチャブルTV』の潜入取材を録画して見せてくれたことだったんだけど。今はもう、帰省するたびに喧嘩になる。生活が後ろ倒しになって社会に適合しないとか、服装が汚いとか。でもこれって大学生の間しか出来ないし。こんなにやってることなのに、わかってもらえないの悲しいなって思うんですよね。依存じゃないのって言われたりもして、親からは悪い印象もたれてますね。

\talkerw{Z} 思想強そうとか言われがち。

\talkerw{ぴっぴき} まあでも思想は全然。全員が全員思想をもっている訳じゃない。

\talkerw{午後3時} 意外と、400人以上住んでて女子もめっちゃいる、っていうのはあんまり知らない人が多いかな。あと、最初から言わずにあとあとさらっと寮住んでるよって言うことによって、「こういう人もいるんだ」っていうのを緩やかに主張しようと思って。

\talkerw{へちま} 私も。たくさんの人間に私は熊野寮に住んでますって言うと、「あっこういう人もいるんだ」みたいな、知ってもらえるのかなーと思っているフシがあります。

\talkerw{午後3時} この座談会はそういう人のためにやってるとも言えるから。パンフレット面白いよ。入寮パンフレットめっちゃ好きで、受験の時に吉田寮も地塩寮も全部で9冊もらって帰った。

\talkerw{みかん} 参考書よりカサをとってる。

\talkerw{午後3時} 座談会読んでた人達と寮入ってから会ったりすると、あー！っていう面白さがある。ぜひ入ってから会ってください、この5人に。

\talkerw{ぴっぴき} 自分はパンフレット見てて、「この人、Twitterでコンタクトとった人と一緒じゃねぇ!?」って思ってたら、実際入ったらそうやったみたいなのがあって。納得だわ！みたいな。笑　

\talkerw{へちま} 住んでから読み返してみると、書きぶりで「これ多分この人だな」って結構わかるようになりました。

\noindent{\uuline{\Large\textbf{家族からの印象\\}}}

\talkerw{Z} 私(熊野に住んでいることを)言ってないです。偏見強過ぎて結構反対派みたいな。親の世代とかは学生運動の影響であまりいいイメージ持たないから。

\talkerw{へちま} うちの親は、安い寮があると思ってる。「安い寮あるよ」「じゃあそこに入りなよ」ってなって入った。それ以降はなんも話題にならない。

\talkerw{みかん} なんとなくいい感じに言ってる。帰省すると「ワイルドになったね」と言われる笑。

\talkerw{ぴっぴき} どういう点ですか？

\talkerw{みかん} なんか服装とか、焼けたり。でも寮はセキュリティとかも、一人暮らしよりも人の目があるからいいですよね。

\talkerw{午後3時} 隣に住んでる人とかめっちゃわかる。

\talkerw{みかん} 具合悪くなっても、隣人が声かけてくれたり、厚生部\footnote{シャワー室管理や寮内の物品補充、寮の衛生をつかさどる部会。熊野寮は部会と委員会、局があって仕事に専門性がある。}棚に薬が入ってたり。

\talkerw{午後3時} 2020年のコロナ流行時に引っ越してすぐ熱が出た時、同部屋の先輩がいろいろ買ってきてくれて。薬局も近いし、先生が優しい病院も教えてくれた。婦人科系のこととかも、女子寮生グループで聞ける人がいたり助けてくれる人がいたりするのがありがたいなと思ってます。

\talkerw{午後3時} うちは親が熊野寮好き。なんなら親が最初に探してきた。親が自分が住みたかったって感じだし、くまのまつりとか呼ぶまでもなく来ます。

\talkerw{ぴっぴき} 休学とか、留年の意識がすごい軽いよね。

\talkerw{午後3時} 休学したいって思ったときに、例がいっぱいあるのはいいこと。寮にいるから留年するというよりは、留年しても何とかなることが分かるという感じがあると思う。

\talkerw{みかん} 5とか6とか聞くと嬉しくなる。その先へ。

\talkerw{へちま} 私の先輩は一時期研究室に行けなくなった時、セーフティーネットとして機能したのは寮だったと言ってました。

\talkerw{みかん} 楽しい寄り道みたいな。バリバリ就活して4回生で卒業って人もいるし、どっちがいいとか規定されない感じがいい。

\talkerw{ぴっぴき} 自分、人生の走馬灯に絶対熊野寮での生活が出てくる。いろんな価値観を知れたり、深い話ができたり。尊敬してる人がすごい多いから。エネルギーがすごいこの寮にはある気がするんだよなあ。

\talkerw{みかん} 無限に頭いい人、議論強い人が多すぎる。サークルと違って無限にいろんな人がいる。

\talkerw{Z} 合わない人がいても、母数が多いし幅が広いから、合う人が必ず出てくるとは思うかな。

\noindent{\uuline{\Large\textbf{寮の暮らしにくさ\\}}}

\talkerw{みかん} トイレが、食堂最寄りが3階だから、ふらっと行くにはちょっと。

\talkerw{へちま} 寮かえる直前に尿意を催したときとかに結構頑張らないと。

\talkerw{午後3時} トイレは去年アンケートとって改善された。前はオールジェンダートイレの個室が並んでて使いにくい話があったけど、今は女子トイレの階を近くして、オールジェンダー側もドアに鍵がかかるようにしてくれた。

\talkerw{Z} 洗濯物関係で。名前書いてあったのに、開けられて下着も全部出されたことがあって。出したのは女性だったんですけど。1時間以内に回収しに行ったのに、丸出しで隣に置かれてたのは、えーってなりました。

\talkerw{ぴっぴき} 自ブロックは炊事場に洗濯機があって男女共用。名前書くのって女子の名前ってわかっちゃうから、それも怖くないですか？

\talkerw{へちま} うちのブロックは「放置されてたら出していい」ってルール。人のを触るのが苦手だとちょっと大変。

\talkerw{ぴっぴき} 最近は、女子シャワー室に小さい洗濯機が設置されたから便利になってるのかな。談話室\footnote{各ブロックに一つあるリビングのようなもの。談話する部屋。}の住みよさとかどうですか？

\talkerw{みかん} ほぼ全員男子のときもあり、女の子のアニメとか苦手な人いるかも。

\talkerw{午後3時} うちのブロックは、テレビじゃなくてスクリーンで上映会みたいにしちゃうから、部屋全体暗くしますみたいな感じになりがち。

\talkerw{へちま} 自分と異なる意見を持つ人の数が多いと、圧倒されて帰ってしまう時はある。

\talkerw{ぴっぴき} 自分は談話室ではワーって叫んだりごろごろしたり、のびのびしてます。でも、男子が多い空間だと、もやもやを言葉にするのにハードルを感じることがたまにある。ちゃんと言葉で理屈で説明しなきゃ、っていう圧が怖いというか。議論強い人が多すぎて、釈明できないのが怖いから、寮の議論というものに対しての怖さがある。

\talkerw{午後3時} 近くに先人がいたらいいなとは思う。同ブロックに場に入っていける女子の先輩がいてくれるのは本当にありがたい。

（〜お菓子を食べながらおすすめの店の話〜）

\noindent{\uuline{\Large\textbf{「ここのラーメンべっぴんさんや！」\\}}}

\begin{itemize}
    \item ぴっぴき：古着屋flamingo／微風台南（雰囲気がいい）
    \item Z：からこの唐揚げ
    \item みかん：ゼスト御池／丸太町パクチー（エスニック）
    \item へちま：山岡家のつけ麺（すっぱい）、賄い御饌（ラーメン、北部キャンパス近く）
    \item 午後３時：洛北阪急スクエア／泉輪、カフェコレクションのバターライス
\end{itemize}

\end{multicols}

\noindent{\uuline{\Large\textbf{最後に\\}}}

この座談会は女子寮生向け相談窓口\footnote{女子寮生窓口＝生活上の悩みやハラスメントの相談受付、女子寮生新歓の開催などを行う団体。}という団体が主体となり開催されました。熊野寮において女子寮生の数は全体の2割程度とマイノリティーであり、女子という属性を持つことで、生活をしていてもやもやした思いを抱くことも少なくありません。女子寮生窓口は匿名でこうした生活上の悩みやハラスメントの相談受付を行ったり、女子寮生新歓を開催することで女子寮生が集まれる場を作ったりと、女子寮生のセーフティーネットとして在れるように活動しています。この座談会が入寮を迷っている女子学生の参考となり、今後の寮生活をイメージする上でお役に立てば幸いです。

以下は個人的に書くことなのですが、自分は寮に入ってよかったと思える最大の恩恵の一つとして、男女の意識にそこまで囚われずにすむようになったというのがあります。寮では原理的に、対等ないち人間としての付き合いや関わりが求められるので、生活をしていてまず性別より先に人間であることを意識することが多いです。集団に女子一人であることを自分も周囲も特になんとも思ってない、というのが寮でよくある風景でしょう。男子大勢のなかに分け入っていくことや、この寮が誰でも思うことがあれば大いに意見を言って自分でも変えていける環境だということがあり、入寮してからというもの自分のなかの強さがきゅっと固まったなあという気がしています。このことは結構私の人生のなかで転換点であって、より一層自由に、自分の人生の主体として人生を彩れるようになったし、目の前の人のことを、その人なりの葛藤やむずかしさを抱えながら生きる同じひとりの人間としてより一層耳を傾けて向き合えるようになったと思っています。いろいろな価値観に揉まれることは時に苦しさであり、思うようにいかないこともきっと経験することでしょうが、ここにはまだ見ぬたくさんのきらめきがあります。入寮を迷っている人はどうかいっぺん入ってみてください。熊野寮であなたとお会いできることを楽しみにしています。

\rightline{ぴっぴき}